\section{Single Qubit States}

Before introducing qubits, we briefly describe what we mean by a \textbf{quantum system}.

\highspace
A \definition{Quantum System} is a \hl{physical system} (such as a photon, an electron, or an atom) \hl{whose behavior is described by the laws of quantum mechanics.}

\highspace
In this course, we will focus on the simplest case: systems that can produce \textbf{two possible outcomes} when measured. These are called \textbf{two-level systems}, and they form the basis of qubits.

\longline

\subsection{Building and Measuring Qubits (Intuition)}

Before introducing the mathematical formalism of qubits, it is useful to develop an \textbf{intuitive understanding} of how a simple quantum system behaves when it is \textbf{prepared} and \textbf{measured}.

\highspace
\begin{flushleft}
    \textcolor{Green3}{\faIcon{question-circle} \textbf{A Simple Quantum Experiment}}
\end{flushleft}
Consider a very simple experimental setup:
\begin{itemize}
    \item A \textbf{source} emits a quantum particle (for example, a photon).
    \item The particle passes through a device.
    \item A \textbf{detector} measures the outcome.
\end{itemize}
Each time we perform the experiment, the detector returns one of two possible results:
\begin{equation*}
    0 \quad \text{or} \quad 1
\end{equation*}
This motivates the idea of a \textbf{two-level system}, which is the simplest type of quantum system.

\highspace
\begin{flushleft}
    \textcolor{Green3}{\faIcon[regular]{lightbulb} \textbf{Key Observation}}
\end{flushleft}
If we repeat the same experiment many times under identical conditions, we observe that:
\begin{itemize}
    \item The outcome is \textbf{not deterministic}.
    \item Instead, it is \textbf{probabilistic}.
\end{itemize}
For example, we might obtain:
\begin{itemize}
    \item Outcome $0$ with probability $p$
    \item Outcome $1$ with probability $1 - p$
\end{itemize}
This is fundamentally different from classical systems, where the outcome is determined in advance.

\highspace
\begin{flushleft}
    \textcolor{Green3}{\faIcon{question-circle} \textbf{What is a Measurement?}}
\end{flushleft}
A \definition{Measurement} is the process that:
\begin{itemize}
    \item Extracts \textbf{classical information} from a quantum system;
    \item Produces a \textbf{definite outcome} (e.g., $0$ or $1$).
\end{itemize}
An important property of quantum measurement is that it \textbf{affects the system itself}. After a measurement is performed, the \textbf{system is no longer in its original configuration}.

\highspace
\begin{flushleft}
    \textcolor{Green3}{\faIcon[regular]{lightbulb} \textbf{From Intuition to Formalization}}
\end{flushleft}
This simple experiment suggests that:
\begin{itemize}
    \item A quantum system can produce \textbf{two possible outcomes}
    \item The outcomes are governed by \textbf{probabilities}
    \item The system must therefore contain some form of \textbf{internal state} that determines these probabilities
\end{itemize}
In the following example, we will introduce a concrete physical example (photon polarization) that allows us to represent this internal state using \textbf{vectors with complex coefficients}.

\begin{examplebox}[: Example of a Two-Level Quantum System]\phantomsection\label{ex:polarization}
    To make the previous discussion more concrete, we introduce a simple physical system that behaves as a two-level quantum system: the \textbf{polarization of a photon}.

    \highspace
    \begin{flushleft}
        \textcolor{Green3}{\faIcon{question-circle} \textbf{What is a Photon?}}
    \end{flushleft}
    A \definition{Photon} is the elementary particle of light. It carries energy and behaves both as a particle and as a wave. In this course, we will use photons as a simple example of a \textbf{quantum system}, since their properties can be easily controlled and measured.

    \highspace
    \begin{flushleft}
        \textcolor{Green3}{\faIcon{question-circle} \textbf{What is Polarization?}}
    \end{flushleft}
    \definition{Polarization} describes the direction in which the electric field of a light wave oscillates. In simple terms, it can be thought of as the \textbf{orientation} of the light wave. For a photon, polarization can be visualized as a direction in space. In particular, we consider two basic polarization directions:
    \begin{itemize}
        \item \textbf{horizontal} ($\rightarrow$)
        \item \textbf{vertical} ($\uparrow$)
    \end{itemize}
    These two directions form a \textbf{two-level system}, since a measurement can distinguish between them (i.e., we can determine whether the photon is horizontally or vertically polarized).

    \highspace
    \begin{flushleft}
        \textcolor{Green3}{\faIcon[regular]{lightbulb} \textbf{Superposition of Polarizations}}
    \end{flushleft}
    A photon is not restricted to being only horizontal or vertical. In general, it can be in a \textbf{combination} of both states:
    \begin{equation*}
        v = a \, \rightarrow + b \, \uparrow
    \end{equation*}
    Where:
    \begin{itemize}
        \item $a$ and $b$ are \textbf{complex coefficients};
        \item They determine the behavior of the system when measured.
    \end{itemize}
    This is the first example of a \textbf{quantum state}: a system described by a linear combination of basic states.

    \highspace
    \begin{flushleft}
        \textcolor{Green3}{\faIcon{question-circle} \textbf{Vector Representation}}
    \end{flushleft}
    We can represent these polarization states using vectors:
    \begin{equation*}
        \rightarrow \, =
        \begin{bmatrix}
            1 \\
            0
        \end{bmatrix}
        \hspace{2em}
        \uparrow \, =
        \begin{bmatrix}
            0 \\
            1
        \end{bmatrix}
    \end{equation*}
    Therefore, a general polarization state becomes:
    \begin{equation*}
        v =
        \begin{bmatrix}
            a \\
            b
        \end{bmatrix}
    \end{equation*}
    This shows that the state of the system can be described as a \textbf{vector with complex components}, exactly as introduced in Dirac notation.

    \highspace
    \begin{flushleft}
        \textcolor{Green3}{\faIcon{book} \textbf{Geometric Interpretation}}
    \end{flushleft}
    The horizontal and vertical polarization states can be interpreted as two \textbf{orthogonal axes} in a two-dimensional space.

    \begin{center}
        \begin{tikzpicture}[>=Stealth]

            % Axes
            \draw[->, thick] (0,0) -- (3,0) node[right] {$\rightarrow$};
            \draw[->, thick] (0,0) -- (0,3) node[above] {$\uparrow$};

            % Origin
            \node[below left] at (0,0) {$0$};

            % Vector
            \draw[thick, blue, ->] (0,0) -- (2,2) node[above right] {$v = a\rightarrow + b\uparrow$};

        \end{tikzpicture}
    \end{center}
    The point where the axes intersect represents the origin, and any polarization state can be represented as a \textbf{vector starting from the origin}.

    \begin{center}
        \begin{tikzpicture}[>=Stealth]
            % Source
            \node[draw, rounded corners] (src) at (0,0) {Photon};

            % Polarizers
            \node[draw, rounded corners] (A) at (3,0) {A ($\rightarrow$)};
            \node[draw, rounded corners] (B) at (6,0) {B ($\uparrow$)};

            % Screen
            \node[draw, rounded corners] (screen) at (9,0) {Screen};

            % Arrows
            \draw[->] (src) -- (A);
            \draw[->] (A) -- (B);
            \draw[->] (B) -- (screen);

        \end{tikzpicture}
    \end{center}

    In this representation:
    \begin{itemize}
        \item $\rightarrow$ and $\uparrow$ form a \textbf{basis}
        \item Any state is a \textbf{linear combination} of these basis states
    \end{itemize}

    \begin{flushleft}
        \textcolor{Green3}{\faIcon[regular]{lightbulb} \textbf{Measurement with Polarizers}}
    \end{flushleft}
    A \definition{Polarizer} is a device that \textbf{filters light based on its polarization}. A polarizer acts like a \textbf{filter} that allows only one polarization direction to pass and blocks the other. For example:
    \begin{itemize}
        \item A polarizer aligned with the horizontal direction ($\rightarrow$) will allow horizontally polarized light to pass and block vertically polarized light.
        \item A polarizer aligned with the vertical direction ($\uparrow$) will allow vertically polarized light to pass and block horizontally polarized light.
    \end{itemize}
    If the photon is a combination:
    \begin{equation*}
        v = a\,\rightarrow + b\,\uparrow
    \end{equation*}
    And we use a:
    \begin{itemize}
        \item Horizontal polarizer ($\rightarrow$), the photon will pass with probability $|a|^2$ and be absorbed with probability $|b|^2$.
        \item Vertical polarizer ($\uparrow$), the photon will pass with probability $|b|^2$ and be absorbed with probability $|a|^2$.
    \end{itemize}
    This shows that:
    \begin{itemize}
        \item Probabilities are determined by the amplitudes (i.e., the coefficients $a$ and $b$).
        \item Measurement \textbf{modifies the state of the system} (e.g., if the photon passes through the horizontal polarizer, it will be in the state $\rightarrow$).
    \end{itemize}
\end{examplebox}